% =============================================================================
% Bloque 5: Hipótesis y Objetivos (slide 11)
% =============================================================================

\section{Hipótesis y objetivos}

% --- Slide 11: H1 + objetivo general + sub-hipótesis + específicos ---
\begin{frame}[c]{Hipótesis y objetivos}
\setbeamercolor{block title}{fg=white,bg=uaiorange}
\begin{block}{H1 — Hipótesis principal}
\small
Si a los ejemplos que escribe la IA les aplicamos un \textbf{control de calidad en el mapa}
(quedarnos con los bien ubicados), el clasificador mejora \textbf{más que con SMOTE},
el método clásico — y la mejora es real, \textbf{no producto del azar}.
\end{block}

\setbeamercolor{block title}{fg=white,bg=uaibluedark}
\begin{block}{Objetivo general}
\small
Desarrollar y validar ese control de calidad de modo que funcione \textbf{mirando solo el mapa}:
sin importar qué IA escribió el texto, y en escenarios con \textbf{pocos ejemplos por clase}.
\end{block}

\end{frame}
